\documentclass[conference]{IEEEtran}
\IEEEoverridecommandlockouts
% SAPEC 2027 — Track 9: Machine Learning & AI in Power Electronics
% South Asia Power Electronics Conference | IIT Bombay | June 3–5, 2027
% Authors: Sai Manoj Patro, Arup Ratan Bhowmick
% BITS Pilani, Hyderabad Campus — Dept. of EEE

\usepackage{cite}
\usepackage{amsmath,amssymb,amsfonts}
\usepackage{algorithmic}
\usepackage{graphicx}
\usepackage{textcomp}
\usepackage{xcolor}
\usepackage{booktabs}   % for \toprule, \midrule, \bottomrule in tables
\usepackage{multirow}
\usepackage{url}

\def\BibTeX{{\rm B\kern-.05em{\sc i\kern-.025em b}\kern-.08em
    T\kern-.1667em\lower.7ex\hbox{E}\kern-.125emX}}

\begin{document}

% -----------------------------------------------------------------------
%  TITLE
% -----------------------------------------------------------------------
\title{Lookback-Optimized Physics-Informed AI Dispatch Architecture
       for PV-BESS Active Power Ramp-Rate Mitigation
       Under IEEE 1547-2018%
\thanks{This work was carried out as part of the Design-Oriented Project (DOP)
at the Department of Electrical and Electronics Engineering, BITS Pilani,
Hyderabad Campus, under the guidance of Prof.\ Arup Ratan Bhowmick.
No external funding to declare.}
}

% -----------------------------------------------------------------------
%  AUTHORS
% -----------------------------------------------------------------------
\author{%
\IEEEauthorblockN{Sai Manoj Patro}
\IEEEauthorblockA{%
  \textit{Dept.\ of Electrical \& Electronics Engg.}\\
  \textit{BITS Pilani, Hyderabad Campus}\\
  Hyderabad, Telangana, India\\
  f20232462@hyderabad.bits-pilani.ac.in}
\and
\IEEEauthorblockN{Arup Ratan Bhowmick}
\IEEEauthorblockA{%
  \textit{Dept.\ of Electrical \& Electronics Engg.}\\
  \textit{BITS Pilani, Hyderabad Campus}\\
  Hyderabad, Telangana, India\\
  arup.ratan@hyderabad.bits-pilani.ac.in}
}

\maketitle

% -----------------------------------------------------------------------
%  ABSTRACT
% -----------------------------------------------------------------------
\begin{abstract}
Active power ramp-rate violations at the point of common coupling (PCC)
represent a critical grid integration challenge for utility-scale photovoltaic
(PV) plants under IEEE~1547-2018. This paper presents a lookback-optimized,
physics-informed machine learning dispatch architecture that enables a PV-BESS
inverter controller to enforce the $\pm10\%$/min ramp-rate limit with greater
than 99.8\% compliance while simultaneously reducing battery energy storage
system (BESS) capital sizing requirements by 32.7\%. A systematic 36-cell
two-dimensional parametric sweep over lookback window
$x \in \{15,30,45,60,90,120\}$~min and dispatch horizon
$t_d \in \{5,10,15,30,45,60\}$~min identifies a global optimum at
$x^{*}=60$~min, $t_d^{*}=5$~min, achieving $R^2=0.8553$ and
RMSE~$=106.59$~W/m$^2$ on an out-of-time holdout (Oct--Dec 2025) using
262,020 records of 1-minute ground pyranometer irradiance from Kuala Lumpur,
Malaysia. Physical innovations include a Sandia dynamic cell-temperature model,
damped cloud-parallax ray-tracing ($\alpha_{\mathrm{eff}}=0.45$), and a causal
online exponentially-weighted moving-average (EWMA) filter that resolves the
Double Penalty Effect inherent in point-MSE optimization. The predictive
dispatch strategy reduces required BESS capacity from 2.60~MWh to 1.75~MWh---a
direct 32.7\% capital expenditure reduction.
\end{abstract}

\begin{IEEEkeywords}
solar irradiance forecasting, PV-BESS dispatch, ramp-rate control,
IEEE~1547-2018, cloud parallax, LightGBM, Double Penalty Effect,
parametric horizon sweep
\end{IEEEkeywords}

% -----------------------------------------------------------------------
%  SECTION I — INTRODUCTION
% -----------------------------------------------------------------------
\section{Introduction}

The rapid proliferation of utility-scale solar PV plants in tropical Southeast
Asia has introduced severe active power variability on distribution and
sub-transmission feeders. Kuala Lumpur's equatorial convective climate generates
irradiance ramp events exceeding 400~W/m$^2$/min during afternoon thunderstorm
episodes---events that propagate directly into MW-scale power injection
transients at the PCC. IEEE~1547-2018 \cite{b2}, the governing standard for
distributed energy resource (DER) interconnection, mandates that any plant
larger than 500~kVA must not cause the PCC ramp rate to exceed
$\pm10\%$ of rated capacity per minute. For a 10~MWp~DC / 8~MW~AC installation
with an inverter loading ratio (ILR) of 1.25, this translates to a hard limit
of $\pm800$~kW/min.

Conventional BESS controllers that react to measured ramp violations after the
fact---termed \emph{reactive dispatch}---are perpetually one sample behind the
irradiance transient. Static persistence forecasters that simply hold the last
measured value forward improve on no-forecast control but cannot anticipate
convective cell evolution. Neither approach can optimally pre-position BESS
state-of-charge (SoC), leading to SoC starvation during sustained ramp
sequences and consequent grid violations or costly inverter clipping
\cite{b7,b8}.

Furthermore, two systematic modeling errors have been overlooked in prior
literature:
\begin{enumerate}
\item \textbf{The Double Penalty Effect:} Point-wise mean squared error (MSE)
      optimization against a pinhole pyranometer is penalized twice---once for
      missed cloud shadow timing and once for the irreducible spatial mismatch
      between a 2--10~km satellite cloud-fraction grid cell and a sub-centimeter
      pyrheliometer aperture. This imposes a fundamental ceiling of
      $R^2 \approx 0.54$ at 1-min resolution, regardless of model sophistication.
\item \textbf{Static parallax bias:} Cloud shadow cast angles computed from
      cloud altitude and solar azimuth/zenith diverge by up to $8^{\circ}$ from
      observed shadow trajectories due to wind-driven cloud advection and
      turbulent boundary-layer parallax. Uncorrected, this bias shifts predicted
      shadow arrival times by $\pm3$--7~min.
\end{enumerate}

The contributions of this work are:
\begin{enumerate}
\item A Sandia dynamic thermal derating model \cite{b1} that reduces
      overestimation of AC power output by 14.5\% at tropical midday
      ($T_{\mathrm{cell}}=61.2^{\circ}$C), directly improving BESS
      charge pre-positioning accuracy.
\item A damped cloud-parallax ray-tracing algorithm ($\alpha_{\mathrm{eff}}=0.45$)
      that eliminates geometric overshoot in shadow arrival prediction.
\item Resolution of the Double Penalty Effect through utility-block target
      aggregation (5-min and 15-min rolling means), unlocking $R^2$ gains of
      0.14--0.22 absolute over 1-min point targets.
\item A 36-cell 2D parametric sweep that identifies the globally optimal
      lookback--horizon pair ($x^{*}=60$~min, $t_d^{*}=5$~min), eliminating
      ad-hoc hyperparameter selection.
\item Demonstration that predictive dispatch reduces BESS sizing from
      2.60~MWh to 1.75~MWh---a 32.7\% reduction in capital expenditure.
\end{enumerate}

% -----------------------------------------------------------------------
%  SECTION II — PHYSICS-INFORMED FORECASTING ARCHITECTURE
% -----------------------------------------------------------------------
\section{Physics-Informed Forecasting Architecture}

\subsection{Sandia Dynamic Cell Temperature Model}

Tropical solar farms operate in ambient temperatures of 32--38$^{\circ}$C with
annual global horizontal irradiance (GHI) exceeding 1,600~kWh/m$^2$/year.
Standard test condition (STC) power ratings at $T_{\mathrm{cell}}=25^{\circ}$C
systematically overestimate AC output. The Sandia model~\cite{b1} represents
module thermal dynamics as a first-order lag:
\begin{equation}
  T_{\mathrm{cell}}(t) = T_{\mathrm{amb}}(t)
    + \mathrm{GHI}(t)\cdot e^{a + b\cdot WS} + \Delta T,
  \label{eq:tcell}
\end{equation}
where $a=-3.47$, $b=-0.0594$ (glass/glass open-rack), $WS$ is wind speed
(m/s), and $\Delta T=3^{\circ}$C is the junction-to-back-surface temperature
differential. AC power output is derated as:
\begin{equation}
  P_{\mathrm{AC}}(t) = P_{\mathrm{STC}}\cdot\frac{\mathrm{GHI}(t)}{1000}
    \cdot\bigl[1 + \gamma_P\bigl(T_{\mathrm{cell}}(t)-25\bigr)\bigr]
    \cdot\eta_{\mathrm{inv}},
  \label{eq:pac}
\end{equation}
with temperature coefficient $\gamma_P=-0.0035$/°C and inverter efficiency
$\eta_{\mathrm{inv}}=0.97$. During Kuala Lumpur midday (13:00--15:00~LT),
$T_{\mathrm{cell}}$ reaches 61.2$^{\circ}$C, producing a thermal derating of
1.45~MW ($-$14.5\%) relative to STC prediction on a 10~MWp farm. The thermal
lag time constant $\tau=7$~min captures the capacitive thermal mass of the
module backsheet and mitigates spurious power surge forecasts during transient
cloud clearings.

\subsection{Damped Cloud-Parallax Ray-Tracing}

Cloud shadows at ground level are geometrically displaced from the satellite
nadir cloud-fraction pixel by the parallax vector:
\begin{equation}
  \vec{d}_{\mathrm{parallax}}
    = H_{\mathrm{cloud}}\cdot\tan(\theta_z)\cdot\hat{e}_{az},
  \label{eq:parallax}
\end{equation}
where $H_{\mathrm{cloud}}$ is cloud base altitude, $\theta_z$ is solar zenith
angle, and $\hat{e}_{az}$ is the unit vector in the solar azimuth direction.
For tropical low cumulus convection ($H_{\mathrm{low}}\approx1.5$~km), the
displacement at $\theta_z=60^{\circ}$ reaches 2.6~km. However, empirical
validation against 262,020 paired pyranometer--satellite records reveals that
raw geometric parallax consistently overshoots observed shadow positions. We
introduce an effective shrinkage factor:
\begin{equation}
  \vec{d}_{\mathrm{eff}}
    = \alpha_{\mathrm{eff}}\cdot\vec{d}_{\mathrm{parallax}},
    \quad \alpha_{\mathrm{eff}} = 0.45,
  \label{eq:deff}
\end{equation}
accounting for turbulent wind shear, partial cloud optical depth, and sub-pixel
cloud fragmentation. The corrected parallax reduces shadow timing error from
$\pm7.1$~min (raw) to $\pm2.3$~min (damped) in cross-validated evaluation
\cite{b5}.

\subsection{Feature Engineering}

The feature set captures multi-timescale cloud dynamics across eight groups:
clearness index ($k_t=\mathrm{GHI}/\mathrm{GHI}_{\mathrm{clear}}$, $\Delta k_t$);
variability indices ($\mathrm{VI}_{30}=\sigma(k_t,30\,\mathrm{min})$,
$\mathrm{VI}_{5}$); lag diversity
($\mathrm{GHI}_{t-1},\ldots,\mathrm{GHI}_{t-x}$ at log-spaced intervals);
satellite cloud-fraction layers (CC\textsubscript{total}, CC\textsubscript{low},
CC\textsubscript{mid}, CC\textsubscript{high}) from Open-Meteo at 1-hr resolution
interpolated to 1-min; cloud-fraction time derivatives; parallax-corrected shadow
arrival lag; thermal state ($T_{\mathrm{cell}}$, $T_{\mathrm{amb}}$, $WS$,
$\dot{T}_{\mathrm{cell}}$); and Fourier time embeddings
$\sin/\cos(2\pi\cdot\mathrm{DoY}/365)$, $\sin/\cos(2\pi\cdot\mathrm{HoD}/1440)$.

\subsection{LightGBM Ensemble and Causal EWMA Filter}

The core forecaster is a LightGBM gradient-boosted ensemble
($n_{\mathrm{est}}=2000$, $\mathrm{lr}=0.03$, max\_depth$=7$,
num\_leaves$=63$)~\cite{b4} trained on a stacking architecture where a
meta-learner combines LightGBM, XGBoost, and CatBoost base learners.
All models use a custom asymmetric Huber loss that penalizes underprediction
of ramp events $2\times$ more than overprediction, reflecting the asymmetric
cost of BESS under-charging before a ramp.

A causal online EWMA filter is applied post-forecast to remove high-frequency
point-MSE noise without introducing non-causal lookahead bias:
\begin{equation}
  \hat{b}_t = 0.85\cdot\hat{b}_{t-1} + 0.15\cdot e_t,
  \label{eq:ewma}
\end{equation}
where $e_t$ is the raw model residual at time $t$. The EWMA correction delivers
an absolute $R^2$ improvement of 0.07--0.11 across all test horizons.

% -----------------------------------------------------------------------
%  SECTION III — 2D PARAMETRIC SWEEP AND RESULTS
% -----------------------------------------------------------------------
\section{2D Parametric Optimal Horizon Sweep and Results}

\subsection{Dataset and Experimental Protocol}

The dataset comprises \textbf{262,020 records} of 1-minute GHI from a
research-grade Kipp~\&~Zonen CM11 pyranometer at Kuala Lumpur
(3.1390$^{\circ}$N, 101.6869$^{\circ}$E), spanning January--December~2025.
The train/validation/test split follows a strict chronological out-of-time
protocol: \textbf{Train}~Jan--Aug~2025 (176,000 records); \textbf{Val}~Sep~2025
(21,600 records, hyperparameter selection only);
\textbf{Test}~Oct--Dec~2025 (64,420 records, all reported metrics).
No temporal leakage is permitted: the forecast at time $t$ uses only features
observable at or before $t-t_d$.

\subsection{36-Cell Parametric Sweep}

The lookback window $x$ controls the depth of recent irradiance history encoded
as lagged features. The dispatch horizon $t_d$ is the target prediction offset
($t+t_d$). To identify the globally optimal $(x^{*},t_d^{*})$ without ad-hoc
selection, we execute a full grid search:
\begin{equation}
  \mathcal{G} = \{x : 15,30,45,60,90,120\}
    \times \{t_d : 5,10,15,30,45,60\}\;\mathrm{min}.
  \label{eq:grid}
\end{equation}
Each of the 36 cells trains a separate model and is evaluated on the Oct--Dec
holdout. The 5-min CAISO Real-Time Dispatch (RTD) block mean is used as the
aggregated target.

Table~\ref{tab:benchmark} summarizes performance at selected grid points.
The 2D heatmap and 3D response surface are shown in Fig.~\ref{fig:heatmap}.

\begin{table}[htbp]
\caption{Forecast Performance Benchmark at Selected $(x,\,t_d)$ Pairs}
\label{tab:benchmark}
\begin{center}
\renewcommand{\arraystretch}{1.15}
\begin{tabular}{lcccc}
\toprule
\textbf{Method} & $x$ & $t_d$ & $R^2$ & RMSE\\
 & (min) & (min) & & (W/m$^2$)\\
\midrule
Persistence & -- & 5  & --     & 133.97\\
Persistence & -- & 60 & --     & 235.55\\
Phase 7c (EWMA) & 60 & 5  & 0.7164 & 149.30\\
Phase 7c (EWMA) & 60 & 60 & 0.7633 & 131.31\\
\midrule
\textbf{Proposed ($x^*\!=\!60$)} & \textbf{60} & \textbf{5}  & \textbf{0.8553} & \textbf{106.59}\\
Proposed ($t_d\!=\!10$) & 90 & 10 & 0.7960 & 126.44\\
Proposed ($t_d\!=\!15$) & 90 & 15 & 0.7587 & 137.39\\
Proposed ($t_d\!=\!30$) & 90 & 30 & 0.7040 & 151.95\\
Proposed ($t_d\!=\!45$) & 60 & 45 & 0.6697 & 160.63\\
Proposed ($t_d\!=\!60$) & 60 & 60 & 0.6445 & 167.15\\
\bottomrule
\end{tabular}
\end{center}
\end{table}

\begin{figure}[htbp]
  \centerline{\includegraphics[width=\columnwidth]{fig_heatmap.png}}
  \caption{2D heatmap of $R^2$ (left) and RMSE (right) over the 36-cell
           $(x,t_d)$ sweep. Global optimum marked at $x^*=60$~min,
           $t_d^*=5$~min ($R^2=0.8553$). The 3D response surfaces (Figs.~79
           and~80 in the project repository) are provided as supplementary
           material.}
  \label{fig:heatmap}
\end{figure}

\subsection{Key Findings}

The 2D response surface reveals three physical relationships:

\textbf{(1) Convective cloud lifetime governs $x^*$:} Tropical KL convective
cell lifetimes follow a Weibull distribution with mode $\approx10$--20~min and
tail to 45--70~min. A lookback of $x=60$~min captures one full convective
lifecycle---sufficient momentum for trajectory extrapolation but not so long
that stale cloud state contaminates the feature vector. Shorter lookbacks
($x=15$~min) miss the buildup phase; longer ones ($x=120$~min) dilute the
signal with historical states that no longer influence the current shadow
position.

\textbf{(2) $R^2$ plateau near $x=60$~min is robust:} Across all $t_d$
values, $R^2$ peaks within the $x\in\{45,60,90\}$~min band and falls off
symmetrically for $x<30$~min and $x>90$~min, confirming that the 60-min
window is a physically meaningful characteristic scale.

\textbf{(3) Skill vs.\ persistence increases monotonically with $t_d$:}
The persistence RMSE baseline grows faster with $t_d$ than model RMSE;
physics-informed features maintain skill at longer horizons where naive
baselines collapse.

\textbf{Global Optimum:} $x^*=60$~min, $t_d^*=5$~min $\Rightarrow$
$R^2=0.8553$, RMSE $=106.59$~W/m$^2$, Skill $=+20.44\%$ vs.\ persistence.

% -----------------------------------------------------------------------
%  SECTION IV — PV-BESS DISPATCH AND IEEE 1547 COMPLIANCE
% -----------------------------------------------------------------------
\section{PV-BESS Predictive Dispatch and IEEE 1547-2018 Compliance}

\subsection{System Configuration}

The reference plant is a 10~MWp~DC / 8~MW~AC utility-scale PV farm: ILR~$=1.25$,
$\eta_{\mathrm{inv}}=97\%$, $\gamma_P=-0.35\%/^{\circ}$C, BESS round-trip
efficiency 92\%, IEEE~1547-2018 ramp limit $\pm800$~kW/min, 1-minute control
cycle.

\subsection{Thermal Derating Impact}

Under Kuala~Lumpur's 34$^{\circ}$C ambient conditions, the Sandia model
predicts midday cell temperatures reaching $T_{\mathrm{cell}}=61.2^{\circ}$C.
Substituting into \eqref{eq:pac}:
\begin{equation}
  P_{\mathrm{AC,derated}} = P_{\mathrm{STC}}
    \cdot[1+(-0.0035)(61.2-25)]
    \cdot0.97 = 8.55\;\mathrm{MW},
  \label{eq:derate}
\end{equation}
yielding a thermal power loss of \textbf{1.45~MW ($-$14.5\%)} relative to STC
prediction. Failure to account for this derating causes the BESS dispatch
controller to anticipate 1.45~MW more generation than arrives, triggering
unnecessary discharge and SoC depletion before the actual cloud-driven ramp
event. Fig.~\ref{fig:thermal} illustrates the derating profile.

\begin{figure}[htbp]
  \centerline{\includegraphics[width=\columnwidth]{fig_thermal.png}}
  \caption{Thermal derating profile and AC power forecast vs.\ measured output
           during a representative Kuala Lumpur midday (13:00--15:00~LT).
           Sandia model ($\tau=7$~min thermal lag) accurately tracks the
           1.45~MW derating peak at $T_{\mathrm{cell}}=61.2^{\circ}$C.}
  \label{fig:thermal}
\end{figure}

\subsection{Three-Controller Comparative Study}

Three BESS dispatch strategies are evaluated on the Oct--Dec 2025 test period:
(i)~\textbf{Uncontrolled}---no BESS, raw PV output at PCC;
(ii)~\textbf{Reactive}---BESS responds to measured ramp violations after they
occur (1-sample lag);
(iii)~\textbf{Predictive (Proposed)}---BESS pre-positions SoC 5~min ahead using
the optimal $(x^*=60, t_d^*=5)$ forecast.

Table~\ref{tab:bess} and Fig.~\ref{fig:pareto} summarize compliance and sizing.

\begin{table}[htbp]
\caption{IEEE 1547-2018 Ramp Compliance and BESS Sizing}
\label{tab:bess}
\begin{center}
\renewcommand{\arraystretch}{1.15}
\begin{tabular}{lccc}
\toprule
\textbf{Controller} & \textbf{Violation} & \textbf{BESS Size} & \textbf{CAPEX}\\
 & \textbf{Rate} & \textbf{(MWh)} & \textbf{Index}\\
\midrule
Uncontrolled           & 14.8\%       & ---  & ---   \\
Reactive               & 3.2\%        & 2.60 & 1.00$\times$ \\
\textbf{Predictive (Proposed)} & \textbf{$<$0.2\%} & \textbf{1.75} & \textbf{0.673$\times$}\\
\bottomrule
\end{tabular}
\end{center}
\end{table}

\begin{figure}[htbp]
  \centerline{\includegraphics[width=\columnwidth]{fig_pareto.png}}
  \caption{BESS sizing Pareto frontier (compliance rate vs.\ required storage
           capacity). The predictive controller achieves $>$99.8\% IEEE~1547
           compliance at 1.75~MWh---32.7\% less than the reactive controller
           at 2.60~MWh. Right inset: BESS dispatch during a convective storm
           ramp event, with predictive smoothing applied 5~min ahead of cloud
           shadow arrival.}
  \label{fig:pareto}
\end{figure}

The predictive controller achieves $>$99.8\% IEEE~1547 compliance while
requiring a BESS rated at only 1.75~MWh---a \textbf{32.7\% reduction} relative
to the 2.60~MWh required by the reactive controller. At market lithium iron
phosphate (LFP) battery pack prices of approximately \$150--180/kWh, the
850~kWh reduction represents \$127,500--\$153,000 in avoided battery cost per
installation.

\subsection{Operational Workflow}

The proposed dispatch architecture operates as a continuous feedback loop on a
60-second control cycle (steps complete within 200~ms):
\begin{enumerate}
  \item Feature acquisition (1-min GHI, satellite cloud fractions, cell
        temperature, ambient conditions).
  \item Physics preprocessing: damped parallax correction, thermal derating,
        clearness index computation.
  \item Forecast ($t\!\to\!t{+}5$~min): LightGBM ensemble with causal EWMA
        bias correction \eqref{eq:ewma}.
  \item BESS command generation: compute charge/discharge power to maintain
        PCC ramp $\leq\pm800$~kW/min over the forecast horizon.
  \item Inverter setpoint dispatch via IEC~61850 GOOSE message to the
        grid-forming inverter.
  \item Residual correction: update EWMA bias filter with realized forecast
        error.
\end{enumerate}

% -----------------------------------------------------------------------
%  SECTION V — CONCLUSION
% -----------------------------------------------------------------------
\section{Conclusion}

This paper has presented a lookback-optimized, physics-informed machine learning
dispatch architecture for active power ramp-rate control in a 10~MW PV-BESS
plant under IEEE~1547-2018. The principal contributions are:
\begin{itemize}
  \item A systematic 36-cell 2D parametric sweep identifying the globally
        optimal lookback--horizon pair ($x^*=60$~min, $t_d^*=5$~min),
        achieving $R^2=0.8553$ and RMSE~$=106.59$~W/m$^2$ on out-of-time
        test data, surpassing persistence by 20.44\%.
  \item Identification and resolution of the Double Penalty Effect through
        utility-block target aggregation, delivering absolute $R^2$
        improvements of 0.14--0.22.
  \item A damped cloud-parallax algorithm ($\alpha_{\mathrm{eff}}=0.45$) and
        Sandia thermal model that together reduce systematic forecast bias
        by 14.5\% at tropical midday.
  \item Demonstration that predictive BESS dispatch achieves $>$99.8\%
        IEEE~1547-2018 ramp compliance with a 32.7\% reduction in required
        BESS capacity (2.60\,$\to$\,1.75~MWh), saving approximately
        \$130,000 per 10~MW installation.
\end{itemize}

Future work includes: (i)~real-time deployment on a TI~C2000 DSP-based
grid-forming inverter for hardware-in-the-loop validation;
(ii)~extension to probabilistic forecasts via conformal prediction intervals
for risk-aware BESS pre-positioning;
(iii)~transfer learning to other tropical sites (Chennai, Singapore, Colombo).

% -----------------------------------------------------------------------
%  ACKNOWLEDGMENT
% -----------------------------------------------------------------------
\section*{Acknowledgment}

The authors thank the BITS Pilani, Hyderabad Campus Department of Electrical
and Electronics Engineering for providing computational resources and research
guidance. Irradiance data were acquired from publicly available station records;
satellite cloud-fraction data were sourced from the Open-Meteo historical API.

% -----------------------------------------------------------------------
%  REFERENCES
% -----------------------------------------------------------------------
\begin{thebibliography}{00}

\bibitem{b1}
D.~L.~King, W.~E.~Boyson, and J.~A.~Kratochvil,
``Photovoltaic array performance model,''
Sandia National Laboratories, Tech.\ Rep.\ SAND2004-3535, 2004.

\bibitem{b2}
IEEE Std 1547-2018,
\emph{Standard for Interconnection and Interoperability of Distributed Energy
Resources with Associated Electric Power Systems Interfaces},
IEEE, New York, NY, USA, 2018.

\bibitem{b3}
T.~Haupt and K.~P.~Schneider,
``Short-term solar irradiance forecasting using satellite imagery,''
\emph{Sol.\ Energy}, vol.~191, pp.~707--719, 2019.

\bibitem{b4}
G.~Ke \emph{et al.},
``LightGBM: A highly efficient gradient boosting decision tree,''
in \emph{Proc.\ Adv.\ Neural Inf.\ Process.\ Syst.\ (NeurIPS)},
2017, pp.~3146--3154.

\bibitem{b5}
E.~Lorenz, J.~Hurka, D.~Heinemann, and H.~G.~Beyer,
``Irradiance forecasting for the power prediction of grid-connected
photovoltaic systems,''
\emph{IEEE J.\ Sel.\ Top.\ Appl.\ Earth Obs.\ Remote Sens.},
vol.~2, no.~1, pp.~2--10, 2009.

\bibitem{b6}
B.~Espinar, J.-L.~Ram\'{i}rez, A.~Drews, H.~G.~Beyer, and L.~F.~Zarzalejo,
``Analysis of different error parameters applied to solar radiation data from
satellite and German radiometric stations,''
\emph{Sol.\ Energy}, vol.~83, pp.~118--125, 2009.

\bibitem{b7}
R.~Ahmed, V.~Sreeram, Y.~Mishra, and M.~D.~Arif,
``A review and evaluation of the state-of-the-art in PV solar power
forecasting: Techniques and optimization,''
\emph{Renew.\ Sustain.\ Energy Rev.}, vol.~124, p.~109792, 2020.

\bibitem{b8}
M.~Pomella, S.~Patel, and A.~Bhowmick,
``Ramp-rate control strategies for utility-scale PV with BESS under
distribution feeder constraints,''
in \emph{Proc.\ IEEE Energy Convers.\ Congr.\ Expo.\ (ECCE)}, 2023.

\end{thebibliography}

\end{document}
